\documentclass[12pt]{article}
\usepackage[T1]{fontenc}
\usepackage[utf8]{inputenc}
\usepackage{lmodern}
\usepackage{textcomp}
\usepackage{enumitem}
\usepackage{geometry}
\usepackage{graphicx}
\usepackage{amsmath, amssymb}
\geometry{margin=1in}
\begin{document}
\begin{center}
Physics - Undergraduate Level Assessment 23 \
Total Marks: 40 \
Answer all questions.
\end{center}

\section*{Multiple Choice (10 marks)}
\begin{enumerate}[label=\arabic*.]
\item The sign of the charge carriers in a doped semiconductor can be deduced by measuring which of the following properties?
\begin{enumerate}[label=(\alph*)]
    \item Magnetic susceptibility
    \item Hall coefficient
    \item Electrical resistivity
    \item Thermal conductivity
\end{enumerate}
\item A grating spectrometer can just barely resolve two wavelengths of 500 nm and 502 nm, respectively. Which of the following gives the resolving power of the spectrometer?
\begin{enumerate}[label=(\alph*)]
    \item 2
    \item 250
    \item 5,000
    \item 10,000
\end{enumerate}
\item Which of the following ions CANNOT be used as a dopant in germanium to make an n-type semiconductor?
\begin{enumerate}[label=(\alph*)]
    \item As
    \item P
    \item Sb
    \item B
\end{enumerate}
\item A three-dimensional harmonic oscillator is in thermal equilibrium with a temperature reservoir at temperature T. The average total energy of the oscillator is
\begin{enumerate}[label=(\alph*)]
    \item (1/2) k T
    \item kT
    \item (3/2) k T
    \item 3 kT
\end{enumerate}
\item An organ pipe, closed at one end and open at the other, is designed to have a fundamental frequency of C (131 Hz). What is the frequency of the next higher harmonic for this pipe?
\begin{enumerate}[label=(\alph*)]
    \item 44 Hz
    \item 196 Hz
    \item 262 Hz
    \item 393 Hz
\end{enumerate}
\end{enumerate}

\section*{True / False (10 marks)}
\textit{Answer True or False}
\begin{enumerate}[label=\arabic*.]
\setcounter{enumi}{5}
\item True or False: A particle moving at v=0.60 c in the lab frame will travel 450 m before decaying.
\item True or False: The eigenvalues of a Hermitian operator can be imaginary in certain cases.
\item True or False: The resolving power of the spectrometer for the given wavelengths is 250.
\item True or False: The energy required to ionize helium by removing one electron is 39.5 eV.
\item True or False: The rest mass of a particle with total energy 5.0 GeV and momentum 4.9 GeV/c is approximately 1.0 GeV/$c^2$.
\end{enumerate}

\section*{Long Form Response (20 marks)}
\textit{Answer each question with a clear and complete explanation.}
\begin{enumerate}[label=\arabic*.]
\setcounter{enumi}{10}
\item A student measures the disintegration of a radioactive isotope over 10 one-second intervals, obtaining values of 3, 0, 2, 1, 2, 4, 0, 1, 2, and 5. Analyze how long the student should count to achieve a measurement uncertainty of 1 percent, and explain the factors influencing this determination.
\item Discuss the advantages of using a dye laser for spectroscopy over a range of visible wavelengths, and analyze its effectiveness compared to other laser types.
\end{enumerate}

\vfill
\noindent\textit{End of Paper}
\end{document}